% !TEX encoding = UTF-8 Unicode
% !TEX spellcheck = en-US


% This is the root file of your thesis: thesis.tex
% A line starting with % is a comment. In some cases, I have included a command preceded by a %. You may activate the command by removing the %.

%%===================================
\documentclass[12pt]{report}
\usepackage[utf8]{inputenc} % This defines the font-encoding you prefer to use
\usepackage{fourier}
\pagestyle{headings}
\usepackage{marvosym}
\usepackage[pdftex]{graphicx}
\usepackage[margin=2.5cm]{geometry}
\usepackage{rotating}
\usepackage{amsmath}
\usepackage{amssymb}
\usepackage{caption}
\usepackage{todonotes}
\usepackage{booktabs}
\usepackage{subcaption}
%\usepackage{footnote}
\usepackage{bm}
\usepackage{gensymb}
\usepackage{multirow}
\usepackage{url}
\usepackage{natbib}
%\usepackage[numbers]{natbib} % Use this command if you want to refer to the literature by numbers, e.g., [1]
\usepackage{setspace}
\linespread{1.5}

\newenvironment{remark}{\par\mbox{}\par\noindent{\bfseries \emph{Remark}\/}:\normalfont}{\par}
\newenvironment{defin}{\par\mbox{}\par\noindent\bfseries \emph{\Pointinghand\ }\normalfont}{\par}
\setcounter{tocdepth}{2} 
\usepackage[colorlinks=true, a4paper=true, pdfstartview=FitV,
linkcolor=blue, citecolor=blue, urlcolor=blue]{hyperref}
\pdfcompresslevel=9
\fancyhf{}
\renewcommand{\dateseparator}{.}
\fancyhead[L]{\today\ \currenttime}
\fancyhead[R]{\thepage}
\pagestyle{fancy}
%%===================================
\begin{document}
\thispagestyle{empty}

\vspace*{3cm}

\begin{center}
	{\LARGE{}CyberShip Arctic Drillship}
	\par\end{center}{\LARGE \par}

\begin{center}
	{\LARGE{}User Manual }
	\par\end{center}{\LARGE \par}

\begin{figure*}[h!]
	\centering
	\includegraphics[width=\linewidth]{fig/CSAD_photo.png}
\end{figure*}
\begin{center}
{\today }
\end{center}
\begin{flushleft}
	\vfill{}
	\par\end{flushleft}

\begin{flushleft}
	\includegraphics[scale=0.6]{fig/NTNU_logo.pdf}
	\par\end{flushleft}
Faculty of Engineering Science and Technology\\
Department of Marine Technology

\clearpage{}

\setcounter{page}{0}
\pagenumbering{roman}
\addcontentsline{toc}{section}{Preface}
\section*{Preface}
The purpose of this document is to provide a manual that ease the process of using CyberShip Arctic Drillship(CSAD), and concerns only software and hardware of CSAD specifically. For information about the Marine Cybernetics Laboratory(MCLab) and how to implement custom control systems on the vessel, the reader is referred to the CSAD ROS repo, which can be found on GitHub: \url{https://github.com/NTNU-MCS/}

\addcontentsline{toc}{section}{Structure of document}
\section*{Structure of document}
This User Manual is divided in three parts: 
\begin{itemize}
	\item Technical description(hardware, software, mathematical models etc.)
	\item Operation manual(launching, operation and demolition instructions)
	\item In the Appendix, a description of the extended IMU system is given(4 IMUs)
\end{itemize}

\addcontentsline{toc}{section}{CSAD main data}
\begin{table*}[htb!]
	\centering
	\caption{CSAD main data}
	\begin{tabular}{ll}
		\hline
		\textbf{Parameter} & \textbf{Value} \\ \hline
		Length over all & 2.578 [m] \\
		Beam & 0.440 [m] \\
		Depth & 0.211 [m]\\
		Design draft & 0.133 [m]\\
		Weight & 127.92 [kg] \\
		Scale & 1:90\\
		RPi IP-address & 192.168.0.123 \\
		Qualisys body\footnotetext{Body-coordinate of highest marker} & (960, -190, -575) [mm]\\ 
		MATLAB Version & 2016b\\
		LabVIEW Version & 2017\\
		VeriStand Version & 2017\\
		\hline
	\end{tabular}
\end{table*}

\addcontentsline{toc}{section}{Lnown errors and further work}
\section*{Known errors and further work}
There are some known errors and weaknesses on CSAD: 
\begin{itemize}
	\item Due to a lack of any indexing feature on the belt-wheels on the servos , the wheels can slide around while the servos are stationary. This means the offset angle can change and the software can loose track of the actual position of the thruster.
	All bolts attaching the belt wheels to the servos have been tightened on the 1.july.2021. However servo nr.4 has loosened only 7 days later, causing it to loose the offset angle.
	\item The power is currently not connected to the power-switch, meaning that the boat will turn on as soon as a battery is connected.
	\item The fixture holding servo 3(front left) is not deep enough, making the screw holding the belt-wheel press against the bottom of the fixture causing the servo to stall. This problem is circumvented by installing two washers in-between the servo and the fixture.
	\item There are new constraints on maximum thrust. From the towing test carried out in Juny 2017, a bollard pull test was performed. It is suggested to implement this new maximum thrust values in the thrust allocation. 
	\item The weight of the vessel is not correct, as it does not take the moon pool into account. 
\end{itemize}
\tableofcontents
\setcounter{page}{0}
\pagenumbering{arabic}
\part{Technical description}
\chapter{Hardware}
\section{Introduction to CSE1}
The CS Enterprise I was initially bought in 2009, as a fully configured model boat named "Aziz" built by Model Slipway. Due to requirements for master and PhD experiments, the model was refitted by \cite{Skaatun2011}. The work performed on CS Enterprise I include, but is not limited to, dynamic positioning systems, maneuvering systems and path following, and navigation with virtual reality.

The vessel is a 1:50 scale model of a tug boat, and is fitted with two Voith Schneider propellers(VSP) astern and one bow thruster(BT). The main dimensions of the vessel are:
\begin{table}[h!]
	\caption{Main dimensions of CSE1}
	\centering
	\begin{tabular}{cc}
		\hline
		LOA & 1.105[m]\\
		B & 0.248 [m]\\
		$\Delta$ & 14.11 [kg]\\\hline
	\end{tabular}
\end{table}

\subsection{Literature}
The development of CSE1 is a product of much research from several theses, which contain complementary information on the theory applied to the system. 
\subsubsection{Journals and conferences}
\begin{itemize}
	\item LOS guidance for towing an iceberg along a straight-line path \citep{OrstenNorgrenSkjetne2014}
\end{itemize}
\subsubsection{Specialization projects and master theses}
\begin{itemize}
	\item Development of a DP system for CS Enterprise I with Voith Schneider thrusters. \citep{Skaatun2011}
	\item Development of a modularized control architecture for CS Enterprise I for path-following based on LOS and maneuvering theory \citep{Tran2013}
	\item Automatic Reliability-based Control of Iceberg Towing in Open Waters \citep{Orsten2014}
	\item Line-Of-Sight-based maneuvering control design, implementation, and experimental testing for the model ship C/S Enterprise I.\citep{Tran2014}
	\item Remote Control and Automatic Path-following for C/S Enterprise I and ROV Neptunus \citep{Sandved2015}
	\item Marine Telepresence System \citep{Valle2015}
	\item Nonlinear Adaptive Motion Control and Model-Error Analysis for Ships-Simulations and MCLab experiments \citep{bjorne2016nonlinear}
	\item Low-Cost Observer and Path-Following Adaptive Autopilot for Ships \citep{mykland2017}
\end{itemize}
\subsubsection{Other}
\begin{itemize}
	\item YouTube video \citep{Skaatun2014}
\end{itemize}

\section{Actuators}
Figure \ref{fig:position_of_actuators} illustrate the position of the actuators, and their distance from CO is given in Table \ref{tab:thruster_distance}. The BT and VSP motor speeds are controlled by an Electronic Speed Control(ESC). The ESC receive their setpoints as pulse-width modulated (PWM) signals from the cRIO digital output module. The VSP blade pitches are controlled by servos. The servos also receive their setpoint as PWM signals.

\begin{minipage}{\textwidth}
	\begin{minipage}[b]{0.49\textwidth}
		\centering
		\includegraphics[width=\linewidth]{fig/actuators_overview.png}
		\captionof{figure}{Position of actuators. Adapted from \cite{Valle2015}}
		\label{fig:position_of_actuators}
	\end{minipage}
	\hfill
	\begin{minipage}[b]{0.49\textwidth}
		\centering
		\begin{tabular}{ccc}
		\hline
		\textbf{Parameter} & \textbf{Symbol} & \textbf{Value}[m]\\\hline
		x length to VSP & $L_{x,VSP}$ & -0.4574\\
		x length to BT & $L_{BT}$ & 0.3875\\
		y length to VSP & $L_{y,VSP}$ & 0.055\\\hline
		\end{tabular}
	\captionof{table}{Position of actuators}
	\label{tab:thruster_distance}
	\end{minipage}
\vspace{0.5cm}
\end{minipage}

\section{Power system}
CSE1 is powered with one 12V 12Ah battery on-board. Some of the components require different voltage, and thus some voltage converters are mounted. However, the setup works as it is, and by connecting the battery to the wires, the whole system is powered. A schematic of the power grid is illustrated in Figure \ref{fig: CSE1 power}, and Figure \ref{fig:battery_connected} show a photo of the battery mounted and connected. 
\begin{figure*}[htb!]
\centering
\begin{subfigure}{.45\linewidth}
	\centering
	\includegraphics[height=0.4\paperheight]{fig/CSE1_power}
	\caption{CSE1 power system}
	\label{fig: CSE1 power}
\end{subfigure}
\begin{subfigure}{0.45\linewidth}
	\centering
	\includegraphics[width=\linewidth]{fig/battery_mounted.jpg}
	\caption{Battery mounted and connected}
	\label{fig:battery_connected}
\end{subfigure}
\caption{Battery system}
\end{figure*}

\section{IMU}
CSE1 is equipped with one Inertial Measurement Unit (IMU) from Analog Devices. The sensor mounted on-board is the ADIS16470 and includes a triaxis gyroscope and triaxis accelerometer. The sensor has built-in compensation for bias, alignment and sensitivity, and thus provides accurate measurements over a temperature range of -10 to +75 degrees Celsius. The sampling rate is set to 100 Hz. The most relevant data is presented in Table \ref{tab:IMU_specifications}, and for supplementary information the reader is referred to the data sheet \cite{adis16470}. The coordinate frame of the sensor is illustrated in Figure \ref{fig:IMU_reference_frame}, with positive directions illustrated by arrows. THe the sensor is mounted with a different orientation than the body frame, as can be seen in Figure \ref{fig:IMU_mounted}. Using the \textit{zyx}-convention, the sensor frame has an orientation relative body frame: $(\phi, \theta, \psi) = (\pi, 0, 0)$. Hence, by using the rotation matrix with these values, the measured accelerations and angular rates can be rotated to the body-frame. 
\begin{table}[htb!]\caption{IMU specifications}\label{tab:IMU_specifications}
	\centering
	\begin{tabular}{c|c|c|c|}
		\cline{2-4}
		& \textbf{Parameter} & \textbf{Typical value} & \textbf{Unit}\\ \cline{1-4}
		\multicolumn{1}{|c|}{\multirow{5}{*}{\textbf{Gyroscopes}}} & Dynamic range & $\pm 2000$ & \degree/sec\\ 
		\multicolumn{1}{|c|}{} & Sensitivity & 10 & LSB/\degree/sec\\ 
		\multicolumn{1}{|c|}{} & Bias stability, $\sigma$ & 8 & \degree/hr\\ 
		\multicolumn{1}{|c|}{} & Angular random walk & 0.34 & \degree/$\sqrt{hr}$\\ 
		\multicolumn{1}{|c|}{} & Output noise & 0.17 & \degree/sec rms\\ \cline{1-4}
		
		\multicolumn{1}{|c|}{\multirow{5}{*}{\textbf{Accelerometers}}} & Dynamic range & $\pm 40$ & g\\ 
		\multicolumn{1}{|c|}{} & Sensitivity & 800 & LSB/g\\ 
		\multicolumn{1}{|c|}{} & Bias stability, $\sigma$ & 0.013 & mg\\ 
		\multicolumn{1}{|c|}{} & Velocity random walk & 0.037 & m/sec/$\sqrt{hr}$\\ 
		\multicolumn{1}{|c|}{} & Output noise & 2.3 & mg rms\\ \cline{1-4}
		
		\multicolumn{1}{|c|}{\multirow{1}{*}{\textbf{Power supply}}} & Operating voltage& $3.3 \pm 0.3$ & V\\ \cline{1-4}
	\end{tabular}
\end{table}
\begin{figure}[htb!]
	\centering
	\begin{subfigure}{0.45\linewidth}
		\centering
		\includegraphics[width=1\linewidth]{fig/IMU_reference_frame.png}
		\caption{IMU reference frame from manufacturer}
		\label{fig:IMU_reference_frame}
	\end{subfigure}
	\begin{subfigure}{0.45\linewidth}
		\centering
		\includegraphics[width=0.9\linewidth]{fig/IMU_mounted.jpg}
		\caption{IMU mounted in the vessel}
		\label{fig:IMU_mounted}
	\end{subfigure}
\caption{Inertial Measurement Unit in CSE1}
\end{figure}

\section{Control system}
The on-board control system consists of the following parts:
\begin{itemize}
	% \item a National Instruments compact reconfigurable input/output (cRIO)
	% embedded controller
	\item a Raspberry Pi 4 (RPi) single-board computer
	\item three electronic speed controllers (ESC)
	\item four servos
\end{itemize}
% \subsection{cRIO}
% The model on-board is the cRIO-9024, and it is connected to 4 FPGA modules for analogue and digital I/O:
% \begin{itemize}
% 	\item NI-9215, used for analog input such as measuring voltage
% 	\item NI-9263, used for reading IMU measurements
% 	\item NI-9401, not used
% 	\item NI-9474, used for sending pwm signal
% \end{itemize}
\subsection{RPi}

The Raspberry Pi provides communication 
with the DS4 controller, the servos and the electronic speed controllers.. It works as an embedded system, and once powered it will start searching for the wireless controller. When connection is established, it will continuously send signals to the servos and ESC based on operation mode. To successfully connect the sixaxis controller to the RPi, wait for the Bluetooth dongle to start blinking before pressing the PS-button on the controller. 

\subsection{ESC}
The ESC's are controlled with PWM signals, based on PWM tick signals. Table \ref{tab:pwm_spec} describe the setup for all ESC's on-board CSE1, and Table \ref{tab:pwm_range} gives the pwm signal range for each ESC.
\begin{table}[h!]
	\centering
	\caption{PWM specification for ESC}
	\label{tab:pwm_spec}
	\begin{tabular}{cccc}
		\hline
		\textbf{Initial value} & \textbf{Scaling} & \textbf{Offset} & \textbf{PWM period} [Ticks] \\ \hline
		0 & 100 & 0 & 800.000\\ \hline
	\end{tabular}
\end{table}

\begin{table}[h!]
	\centering
	\caption{PWM ranges for ESC}
	\label{tab:pwm_range}
	\begin{tabular}{cccc}
		\hline
		& \textbf{ESC\_BT}[\%] & \textbf{ESC\_VSP1}[\%] & \textbf{ESC\_VSP2}[\%]\\ \hline
		\textbf{min} & 7.00 & 3.12 & 3.12\\
		\textbf{neutral} & 7.55 & 5.01 & 5.01\\
		\textbf{max} & 8.10 & 6.90 & 6.90\\ \hline
	\end{tabular}
\end{table}

\chapter{Software}
\section{Introduction}

The control system of CSE1 is built around a framework called Robot Operating System. ROS consists of tools, libraries, and conventions that help in building robot applications. This chapter gives a description of vessels software hierarchy. Note that most of this software is ready to use, and alterations in the software described here should not be necessary, unless specified.


\section{Driver setup}

\subsection{Actuator drivers}
To activate the actuators on the vessel, first SSH to the onboard RPi 4 and launch nodelet manager and the CSE actuator driver:
\begin{lstlisting}[language=bash,basicstyle=\mlttfamily, breaklines=true]
  $ rosrun nodelet nodelet manager __name:=nodelet_manager
  $ rosrun nodelet nodelet load cse_actuator_driver/cse_actuator_driver_nodelet  
  nodelet_manager __name:=cse_actuator_driver
\end{lstlisting}
You should hear a beeping sound if the nodes are succesfully run. 

\subsection{Qualisys drivers}

Next is to launch the Qualisys drivers. These can be found in the Github repository for the MC-Lab. After downloading the package, you can choose to run it from a external computer or on the onboard RPi. To activate it source your workspace, and run the command


\begin{lstlisting}[language=bash,basicstyle=\mlttfamily, breaklines=true]
  $ roslaunch mocap_qualisys qualisys.launch
\end{lstlisting}

\subsection{IMU-drivers}

To read the data from the IMU, their drivers first have to be read. These are written in C++, but are compatible with a ROS-system written in Python or Matlab. Put the \lstinline{imu_driver} package inside your workspace \textbf{src}-directory and run \lstinline{catkin_make}. After the package is built, first source the bash file and then simply run the package.

\begin{lstlisting}[language=bash,basicstyle=\mlttfamily, breaklines=true]
  $ source devel/setup.bash
  $ rosrun imu_driver imu_node
\end{lstlisting}

The driver publishes the measurements to the topic \textbf{/imu}. 


\section{ROS with Python}

A number of base python-packages have been made to help build control-system for the CSE1. These can all be found on the \href{https://github.com/NTNU-MCS}{NTNU-MCS} Github repository. Simply download the packages you needed for your own project. This section will provide an outline of different packages, the overall architecture of a ROS workspace as well as the basic knowledge required to run a ROS control system. 

\subsection{Prerequisites and recommended software}

\textbf{Required:}
\begin{itemize}
    \item Python 3.8.3 +
\end{itemize}
\textbf{Recommended:}
\begin{itemize}
    \item Visual Studio Code with the \textbf{Remote SSH} extension. This enables you to SSH directly in to raspberry via VS Code, allowing you to edit files on the pi in the text-editor.
    \item A computer running linux. This is not required as most software will run on the Raspberry Pi, but can be beneficial for testing and visualizing.
\end{itemize}

\subsection{Sourcing your ROS-distribution}

Before you do anything with ROS, you need to source the installation. On the RP this is done automatically, but in case you are working on your own computer the following steps are necessary: 

\begin{lstlisting}[language=bash,basicstyle=\mlttfamily, breaklines=true]
  $ source /opt/ros/DISTRIBUTION/setup.bash
\end{lstlisting}

\subsection{Building a catkin workspace}

To run ROS-packages we must first define a worspace for them on the Pi. SSH onto the pi and run 

\begin{lstlisting}[language=bash,basicstyle=\mlttfamily, breaklines=true]
  $ mkdir -p workspace/src; cd workspace
  $ catkin_make
\end{lstlisting}

This should build your catkin workspace, and leave you with three sub-directories; \textbf{build}, \textbf{devel} and \textbf{src}. All ROS-packeges are placed in the \textbf{src}-directory. Everytime you add a package to the workspace, rebuild it using \lstinline{catkin_make}.

\subsection{Running packages}

To run your packages, first make sure you have sourced the \lstinline{setup.bash} file:

\begin{lstlisting}[language=bash,basicstyle=\mlttfamily, breaklines=true]
  $ source devel/setup.bash
\end{lstlisting}

Then, to run the package use

\begin{lstlisting}[language=bash,basicstyle=\mlttfamily, breaklines=true]
  $ rosrun <package-name> <node-script>.py
\end{lstlisting}

where <package-name> is the name of the package you want to run, e.g "controller", and the <node-script>.py is the python script that initialisez the rosnode, e.g "controller.py". 

\subsection{Template workspace}

A template workspace consisting of several packages is provided to make developing control systems simpler. The workspace, denoted \textbf{ws\_templates} in the Git-repository, consists of a total of five ROS-packages corresponding two traditional modules in a maneuvering DP-system. These are named:

\begin{itemize}
    \item controller
    \item guidance
    \item observer
    \item thrust\_allocation
    \item simulator 
\end{itemize}

Additionally, there are three packages that provide supporting tools and conventions: 

\begin{itemize}
    \item gain\_server
    \item messages
    \item common\_tools
\end{itemize}

\subsection{Node-structure}

The first four nodes; \textit{controller}, \textit{guidance}, \textit{observer} and \textit{thrust\_allocation} are generic, skeleton packages. Their purpose is to provide a base to build a control-system on, so they only contain the necessary functionality to execute ROS-nodes. The main functionality is left to the user to implement. All the packages follow the exact same structure.

Within every package there is a \textit{src}-directory where a single python script is placed. This is where you will write and call your code. Listing \ref{listing:node} shows an example of a ROS-node that maps the inputs from a PS4 controller to forces in each of CSE1 thrusters. 


\begin{listing}[h]
\inputminted{python}{feedback_ctrl.py}
\caption{Example of a node-file}
\label{listing:node}
\end{listing}

Quickly summarized, you write your functions above the 

\begin{lstlisting}[language=python]
    if__name__ == 'main'
\end{lstlisting}

statement. These functions should then be called in the while loop located inside. The two statements before the while loop are functions for initiating the node, and setting the refresh rate, or frequency, that we want to run the node at. This is predefined in all templates so that you only have to worry about writing the functions and calling them inside the while loop.

Fully fleshed out packages are also available on Github, if one just needs a working control system for experiments without any implementation. These can be found in the folders of \textit{DP\_Lab} and \textit{DP\_Nom} and consist of a backstepping manuevering regulator, a DP-observer, thrust allocation schemes as well as a path-parameterization guidance function. 

\subsection{Simulator}

The simulator package contains a full Model-in-the-loop ROS implementation of the CSE1. It is based of the low-speed dynamics, and can be used for testing the implemented control systems outside of the MC-lab.

The simulator subscribes tp feedback from the topic \textbf{CSEI/u}, to model the signal flow on board the vessel itself. Upon receiving a signal from the topic, it will compute the generalized forces and new position of the vessel.

Note that the model is very simple, and does not take in to account  environmental forces such as current or wind. 


\subsection{Common\_tools package}

As seen in Listing \ref{listing:node}, the templates import from a module called \textbf{common\_tools}. This is where you can find all the tools for publishing and subscribing to the relevant signals in your system, as well as the functions that initialize nodes. The tools are object-oriented, with every relevant signal pertaining to an object in the library. Each object contains methods called publish, and callback. 

Every time a subscriber detects that a new signal is published to its topic, it calls the callback function in our  object, and updates its class variables. You can then retrieve these by using a get-method, or just calling the variable directly through the object.

\begin{lstlisting}[language=python]
  u = u_data.get_data() # Retrieve actuator commands 
  u = u_data.Udata      # Retrieve acutator commands
\end{lstlisting}

The publish method takes an input, and publishes it to the relevant topic.


\begin{lstlisting}[language=python]

 u_data.publish(u) # Publish actuator commands 
\end{lstlisting}

All objects are declared and imported into the relevant code-templates from the \textbf{lib} submodule . Subscriptions to the different topics are all handled inside these classes and the initialization nodes.

It is important that you at any time keep the common\_tools package inside your workspace. The package also contains a submodule called \textbf{math\_tools}. This contains some predefined mathematical functions that may be useful for the project. Rotation matrices and conversions between euler-angles and quaternions are examples of tool you can find here. Feel free to add more functions here as you please, and import them to your nodes. This makes for cleaner and more structured code. 

\subsection{Messages}

While ROS provides many different messaging formats, some times these may be unsuitable to the needs of the system. One might, for example, wish to store the 9 estimated states of the DP observer in a better format than a generic array. That is where custom messages enter. Custom message templates are all stored in the message-package, and can be imported to any other module in the workspace. 

\subsection{The gain\_server package}

To make tuning your controller and observer more effective, tools for dynamic tuning are provided in the form of the ROS-package gain\_server. This allows you to change the gains in real time, rather than having to hard-code the gains and then stopping and starting your simulation every time you want to change them. To activate the node, and GUI use the following steps: 

\begin{itemize}
    \item Make sure you have imported the Gains object from lib.py. It is built like most other objects and contais get-methods for extracting gains, along with callback-functions.
    \item In the \textbf{observerInitNode()} or \textbf{controllerInitNode} uncomment the line starting with \textit{gain\_client}.
    \item To run the node use the command:
    \begin{lstlisting}[language=bash,basicstyle=\mlttfamily, breaklines=true]
    $ rosrun gain_server server.py
    \end{lstlisting}
    \textbf{NOTE!} If you use the dynamic tuning, you should always activate this node \textbf{first} or you may experience errors. This is because the observer and controller nodes will expect values to be there when they are not. When the node is activated, the initial gains should be printed to your terminal
    \item In another terminal, activate the GUI: 
     \begin{lstlisting}[language=bash,basicstyle=\mlttfamily, breaklines=true]
    $ rosrun rqt_gui rqt_gui -s reconfigure
    \end{lstlisting}
    
    This should open a program that looks like this: 
    
    \begin{figure}[h]
        \centering
        \includegraphics[width=\textwidth]{fig/dynamic-reconfigure(1).png}
        \caption{Tuning GUI}
        \label{fig:tuner}
    \end{figure}
    
    Each field or slider here corresponds to a gain in either the observer or controller. Not all the variables are relevant either. The last to are single float values, while the first six are meant to represent the elements on the diagonal of a $3\times3$-matrix. Change these to tune you controller. 
\end{itemize}

\subsubsection{Parameters}

Currently, the package supports a fixed set of parameters. These are observer injection gains $L1, L2, L3$ as well as the controller gains $Kp, Kd, Ki$, $\mu$ and $Uref$. If other parameters are needed these must be added to the \textbf{gains.cfg} file in the cfg-directory. 

\subsection{Launch files}

As a project progresses, more and more nodes are typically added. The process of activating nodes can therefore become more tedious as the complexity increases. To avoid having to open new command-line windows for every node, we can instead use \textit{launch-files} to activate multiple nodes simultaneously. A launch file is easy to create, and is placed in the \textit{launch} directory of a package.

To run a launch file we use the following command: 

\begin{lstlisting}[language=bash,basicstyle=\mlttfamily, breaklines=true]
  $ roslaunch <package_name> my_launchfile.launch  
\end{lstlisting} 

Replace package\_name with the name of the package you placed your launch file in. All launch files have the .launch ending (e.g DP-system.launch). 

\subsection{Usefull ROS-commands}

The following command-line tools are usefull when running a ros system: 

\begin{enumerate}
    \item \begin{lstlisting}[language=bash,basicstyle=\mlttfamily, breaklines=true]
  $ rostopic list \end{lstlisting} - Shows a list of all ROS-topics 
  
    \item \begin{lstlisting}[language=bash,basicstyle=\mlttfamily, breaklines=true]
  $ rosnode list\end{lstlisting} -
    Shows a list of all active ROS-nodes
     \item \begin{lstlisting}[language=bash,basicstyle=\mlttfamily, breaklines=true]
  $ rostopic echo topic\end{lstlisting} - Prints the messages and data sent to each topic to the command-line window
\end{enumerate}



\part{User Manual}\label{part2}
It is assumed that the reader has studied the MCLab Handbook before using CSAD, and has knowledge about Lab equipment, procedures and Safety precautions. In addition, the following is important to keep in mind when using CSAD:
\begin{description}
	\item [{Water~damage:}] CSAD is watertight when the hatches on the top are closed properly.  
	\item [{Propeller~dry~running:}] The thruster gears are lubricated with water, and thus the propellers always has to be in the water when running. Hence, always keep the vessel in the basin when the power is connected. 
	\item [{Total~loss~of~control:}] Pull CSAD with a boat hook, and keep the vessel in water while disconnecting batteries.
	\item [{Launching}] CSAD is a large model, and care must be taken when launching the vessel to the basin. Always be two persons, and make sure the vessel does not hit the basin wall when launching or removing it from the basin. When launching, remove all weights(batteries and ballast). Still, the vessel is heavy, and the lifting up and down to the basin might be harder than expected. 
\end{description}
\chapter{Launching}

\section{Vessel and lab preparations}
Follow these vessel-specific instructions when preparing for experiments with CSAD: 
\begin{enumerate}
	\item Make sure all batteries and weights are removed from the vessel when lifting it. Launch the vessel in the basin. 
	\item Place all 6 batteries(12V 12Ah, marked CSAD) in the vessel, at their dedicated places. 3 in front of the moonpool, 3 behind it. Connect the batteries, positive/7red first then negative/black. 
	\item Place the ballast weights. 20 kg in the aft, and 27.5 in the front. Manually adjust their position, such that the vessel does not have any heel (slagside). Check with the design draft indicated on the outside of the hull. 
	\item run /launch necessary rosnodes for controlling the ship.( The setMotorPower method should be called at least once with 0 as power for each motor.) All motors should beep one note followed by three rising notes, indicating that they are ready to control the motors.
	\item Press the PS-button on the ps4 controller and wait until the indicator light on the controller stops blinking and turns a solid purple colour.
	\item Place the vessel inside the region of sight for Qualisys (check on the Qualisys computer that all 4 reflectors are visible for all cameras). Align the vessel with 0\degree heading in the basin frame, i.e. with the bow pointing towards the command center. 
	\item On the Qualisys computer, aqcuire the body. This process is described in MC-Lab Handbook, with information on debugging. In the body frame, the highest marker has position $(x,y,z)=(960,-190,-575)[mm]$. 
	\item Go to 3D Visualization in QTM, and verify that the body is correct. The body x-axis should be parallel to the lines between the markers on starboard and port side.
\end{enumerate}
CSAD and the lab is now set up for experiments. 

\chapter{Demolition}
When the experiments are finished, follow the procedure given here to shut down. 
\begin{enumerate}
	\item Navigate CSAD near the basin wall
	\item Terminate all rosnodes
	\item Turn of the power switch, and disconnect the batteries
	\item Remove the ballast weights
	\item Remove the batteries from the vessel
	\item Lift CSAD from the basin, and put it in its rack in the storage
	\item Leave the ds4 controller in the vessel
	\item On the Qualisys computer, quit Qualisys Track Manager
	\item If you recorded any videos with the Camera System, export these videos to a memory stick, quit the software and turn of the TV-monitor
	\item Do a general clean up, bring all your personal belongings with you when you leave
\end{enumerate}
NOTE on charging the batteries: When charging the batteries in CSAD, the WiFi bridge must be disconnected. Unplug the power wire from the WiFi bridge(connected on the side of the Ethernet cable). All batteries must be placed in the vessel when charing, and connected, but with the power switch turned off. Connect the charger to the charging wire located in the aft hatch. The charger is located in the shelf in the storage, marked CSAD. Set the charging mode to the motorcycle symbol. 

%%=========================================
\appendix
\appendix
\chapter{Advanced software topics}
\section{Creating FPGA and XML files}
\subsection{Create FPGA target and XML}\label{subsec:Create-FPGA-target}
If you do not have a Veristand FPGA target at your disposal, follow the steps below. If you have a target available and just need to install it in NI Veristand, please jump to Section \ref{sec:install_FPGA_in_veristand}. For CS Enterprise 1 and CS Arctic Drill Ship, the FPGA targets are found on GitHub. 
\begin{enumerate}
	\item Open LabVIEW and create new project. In this guide, LabVIEW 2013 is used, but the procedure should be similar on newer versions. 
	\item Choose NI Veristand FPGA Project in project templates and proceed.
	\item Choose CompactRIO Reconfiguarble Embedded System and click next.
	\item You will now get the choice between letting LabVIEW detect your cRIO  system or configure it yourself. If you are connected to the cRIO and it has all of the I/O ports connected, the option ``Discover existing system'' is simpler and therefore recommneded. If you do not have your cRIO connected choose ``Create new system'', this is the version that will be worked through here.
	\item Select your controller, in our case cRIO-9024.
	\item Select your FPGA target, in our case cRIO-9113.
	\item Then you select your I/O modules to the correct slots. In our case NI 9215 in slot 1 and NI 9474 in slot 4.
	\item You are now finished with configuring your project. Press next.
	\item The project menu will now appear and should look something like Figure
	\ref{fig: Create Labview FPGA target and XML-10}. Select the LabVIEW VI as demonstrated our is called Custom Personality FPGA.vi
	\item The UI window will now present itself, select window and show block diagram.
	\item You should now see a block diagram similar to Figure \ref{fig: Create Labview FPGA target and XML-11}. You will now have to redesign this to look like Figure \ref{fig: Create Labview FPGA target and XML-12}. This will be valid for our system, if you have different I/O modules the block diagram need to reflect this.
	\item Now, return to the Project explorer and select Build Specifications and Custom Personality FPGA
	\item A new window will open. Check that the name and project path is correct and press build.
	\item Select your preferred compile server. The compilation process will take quite some time (approx 15-30 min).
	\item When the compilation process is finished, the last step is to edit the automatically generated XML file. You will now have to find you project directory in Windows. Here there will be a folder called bitfiles which contains the files you compiled in the last step, there will also be a .XML file. The point of editing this file is to match the actually compiled VI, meaning the packets must match the connected I/O. The recommended way to edit the file is to copy our XML file from: Dropbox\textbackslash{}TMR4243 - LAB\textbackslash{}04 cRIO software\textbackslash{}FPGA IO. You will have to make sure that the name of your bitfile matches the name in the XML file as seen in Figure \ref{fig: Create Labview FPGA target and XML-17}, also make sure the I/O modules matches your setup.
	\item Copy the bitfiles from the bitfile folder to the level above so that the bitfile aand the XML file is in the same folder.
\end{enumerate}
Documentation: \url{https://decibel.ni.com/content/docs/DOC-13815}

\begin{figure}[htb!]
	\centering \includegraphics[scale=0.45]{Screenshots/Screenshot_2015-01-16_19-21-16.png}
	\caption{Create Labview FPGA target and XML - 1}
	\label{fig: Create Labview FPGA target and XML-1} 
\end{figure}
\begin{figure}[htb!]
	\centering \includegraphics[scale=0.45]{Screenshots/Screenshot_2015-01-16_19-23-23.png}
	\caption{Create Labview FPGA target and XML - 2}
	\label{fig: Create Labview FPGA target and XML-2} 
\end{figure}
\begin{figure}[htb!]
	\centering \includegraphics[scale=0.45]{Screenshots/Screenshot_2015-01-16_19-23-41.png}
	\caption{Create Labview FPGA target and XML - 3}
	\label{fig: Create Labview FPGA target and XML-3} 
\end{figure}
\begin{figure}[htb!]
	\centering \includegraphics[scale=0.45]{Screenshots/Screenshot_2015-01-16_19-23-58.png}
	\caption{Create Labview FPGA target and XML -4}
	\label{fig: Create Labview FPGA target and XML-4} 
\end{figure}
\begin{figure}[htb!]
	\centering \includegraphics[scale=0.45]{Screenshots/Screenshot_2015-01-16_19-24-31.png}
	\caption{Create Labview FPGA target and XML - 5}
	\label{fig: Create Labview FPGA target and XML-5} 
\end{figure}
\begin{figure}[htb!]
	\centering \includegraphics[scale=0.45]{Screenshots/Screenshot_2015-01-16_19-24-43.png}
	\caption{Create Labview FPGA target and XML - 6}
	\label{fig: Create Labview FPGA target and XML-6} 
\end{figure}
\begin{figure}[htb!]
	\centering \includegraphics[scale=0.45]{Screenshots/Screenshot_2015-01-16_19-25-37.png}
	\caption{Create Labview FPGA target and XML - 7}
	\label{fig: Create Labview FPGA target and XML-7} 
\end{figure}
\begin{figure}[htb!]
	\centering \includegraphics[scale=0.45]{Screenshots/Screenshot_2015-01-16_19-25-54.png}
	\caption{Create Labview FPGA target and XML - 8}
	\label{fig: Create Labview FPGA target and XML-8} 
\end{figure}
\begin{figure}[htb!]
	\centering \includegraphics[scale=0.45]{Screenshots/Screenshot_2015-01-16_19-28-17.png}
	\caption{Create Labview FPGA target and XML -9}
	\label{fig: Create Labview FPGA target and XML-9} 
\end{figure}
\begin{figure}[htb!]
	\centering \includegraphics[scale=0.45]{Screenshots/Screenshot_2015-01-16_19-28-41.png}
	\caption{Create Labview FPGA target and XML - 10}
	\label{fig: Create Labview FPGA target and XML-10} 
\end{figure}
\begin{figure}[htb!]
	\centering \includegraphics[angle=-90,scale=0.45]{Screenshots/Screenshot_2015-01-16_19-29-04.png}
	\caption{Create Labview FPGA target and XML - 11}
	\label{fig: Create Labview FPGA target and XML-11} 
\end{figure}
\begin{figure}[htb!]
	\centering \includegraphics[angle=-90,scale=0.45]{Screenshots/Screenshot_2015-01-16_20-07-43.png}
	\caption{Create Labview FPGA target and XML - 12}
	\label{fig: Create Labview FPGA target and XML-12} 
\end{figure}
\begin{figure}[htb!]
	\centering \includegraphics[scale=0.45]{Screenshots/Screenshot_2015-01-16_19-52-25.png}
	\caption{Create Labview FPGA target and XML - 13}
	\label{fig: Create Labview FPGA target and XML-13} 
\end{figure}
\begin{figure}[htb!]
	\centering \includegraphics[scale=0.45]{Screenshots/Screenshot_2015-01-16_19-53-01.png}
	\caption{Create Labview FPGA target and XML - 14}
	\label{fig: Create Labview FPGA target and XML-14} 
\end{figure}
\begin{figure}[htb!]
	\centering \includegraphics[scale=0.45]{Screenshots/Screenshot_2015-01-16_19-53-25.png}
	\caption{Create Labview FPGA target and XML - 15}
	\label{fig: Create Labview FPGA target and XML-15} 
\end{figure}
\begin{figure}[htb!]
	\centering \includegraphics[scale=0.45]{Screenshots/Screenshot_2015-01-16_20-01-32.png}
	\caption{Create Labview FPGA target and XML - 16}
	\label{fig: Create Labview FPGA target and XML-16} 
\end{figure}
\begin{figure}[htb!]
	\centering \includegraphics[scale=0.45]{Screenshots/Screenshot_2015-01-17_13-59-31.png}
	\caption{Create Labview FPGA target and XML - 17}
	\label{fig: Create Labview FPGA target and XML-17} 
\end{figure}

\subsection{Install in VeriStand}\label{sec:install_FPGA_in_veristand}

The Veristand software does not recognize the physical I/O components of the cRIO. It is necessarry to write a specific FPGA mapping for the specific setup. This results in a XML file that maps the ports. 

To add this file to your Veristand project, enter the system explorer and find the FPGA pane under \textit{targets\textbackslash{}controller\textbackslash{}hardware\textbackslash{}chassis}, as seen in Figure \ref{fig: fpga1}. 
\begin{figure}[htb!]
	\centering \includegraphics[scale=0.45]{fig/fpga1}
	\caption{FPGA1}
	\label{fig: fpga1}
\end{figure}
The next step is to find your XML file. In this case called cRIO-9113 Ex, it is very important that the XML file is placed on level above the FPGA bitfile folder in the directory system, as the files are really being used are the FPGA bitfiles.
\begin{figure}[htb!]
	\centering \includegraphics[scale=0.45]{fig/fpga2}
	\caption{FPGA2}
	\label{fig: fpga2}
\end{figure}
The menu in should now look something like Figure \ref{fig: fpga2}, here you can see the analogue input signals and the digital output PWM signals. These can again be linked to other signals as seen in Figure\ref{fig: veristand mappings}.
\begin{figure}
	\includegraphics[scale=0.45]{fig/fpga3}
	\caption{FPGA3}
	\label{fig: fpga3}
\end{figure}
\subsubsection{Ticks}
tick = FPGA clock pulse
\[\text{tick in seconds}=\frac{1}{\text{frequency}}=\frac{1}{40MHz}=\frac{1}{40*10^{6}}=25*10^{-}9=25ns\]
output at 50 Hz demands output every 
\[\frac{40MHz}{50Hz}=\frac{40*10^{6}}{50}=800000tick\]
\section{Custom Device}
\subsection{Install}\label{subsec: Installing custom device driver}
As of July 2017, there are 3 different Custom Device drivers developed for use in MCLab: 
\begin{description}
	\item [WL\_Joystick] - used for reading sixaxis data sent from the RPi. Created by Torgeir Wahl
	\item [Oqus] - used for reading position and orientation data sent from the Qualisys system. Created by Torgeir Wahl
	\item [IMU] - used for reading IMU data from 4 IMU's, mainly intended for CSAD. Created by Guttorm Udjus
\end{description}
To install a Custom Device driver, the first step is to copy the folder with the driver to the path: \path{C:\Users\Public\Documents\National Instruments\NI VeriStand 201X\Custom Devices}
The directory should now contain something like Figure \ref{fig: custom device folder}.
\begin{figure}[htb!]
	\centering \includegraphics[scale=0.45]{fig/custom_devices_folder.png}
	\caption{Custom device folder}
	\label{fig: custom device folder} 
\end{figure}
The next step is to add custom device to your project. This is done in the system explorer, which is found as seen in Figure \ref{fig: veristand launch system explorer}.
\begin{figure}[htb!]
	\centering \includegraphics[scale=0.45]{fig/"veristand launch system explorer".png}
	\caption{VeriStand launch system explorer}
	\label{fig: veristand launch system explorer} 
\end{figure}
When in the system explorer, adding the custom device should be as simple as right clicking the custom device pane and choosing WL\_Joystick, as in Figure \ref{fig: custom device selection}. If you do not find the custom device WL\_Joystick, the most likely problem is that the placement of the custom device folder from step 1 is wrong.
\begin{figure}[htb!]
	\centering \includegraphics[scale=0.45]{fig/"veristand custom driver select".png}
	\caption{Custom device selection}
	\label{fig: custom device selection} 
\end{figure}
If the installation is successful you should be able to see WL\_Joystick folder under custom devices as seen in the red box in Figure \ref{fig: veristand confirmation}. Here you will also see the different inputs from the custom device, in this case it is joystick axis.
\begin{figure}[htb!]
	\centering \includegraphics[scale=0.45]{fig/"veristand mapping".png}
	\caption{VeriStand}
	\label{fig: veristand confirmation} 
\end{figure}
To connect the joystick to the input ports of the Simulink model. You open the system configuration mappings (click the button marked by the arrow in Figure \ref{fig: veristand confirmation}).
\begin{figure}[htb!]
	\centering \includegraphics[scale=0.45]{fig/"veristand mapping2".png}
	\caption{VeriStand System Configuration Mappings}
	\label{fig: veristand mappings} 
\end{figure}
You the simply find the ports you would like to connect, mark them and click the connect button. Figure \ref{fig: veristand mappings} a joystick output is connected to a input port on the Simulink model.
\section{Raspberry Pi}\label{subsec: RPi setup}
The unit is configured with Raspbian Linux-kernel-based operating system 
\subsection{Raspbian installation and setup}
This section describes how to install and access the Raspbian operating system on the RPi from a Windows computer. The operations are also possible from an OSX or Linux computer.
\subsubsection{Download operating system and utilities}
Download and extract the newest Raspbian\footnote{raspberrypi.org/downloads} operating system (OS) image. Necessary utilities for the setup are
\begin{itemize}
	\item Win32 Disk Imager\footnote{sourceforge.net/projects/win32diskimager} to write the OS image to the RPi SD card
	\item Advanced IP scanner\footnote{by Famatech, advanced-ip-scanner.com} to find the RPi address on the network 
	\item Putty terminal emulator\footnote{www.chiark.greenend.org.uk/\textasciitilde{}sgtatham/putty/download.html} for SSH connection
	\item WinSCP\footnote{by Martin Prikryl, winscp.net/eng/download.php} for file transfer
\end{itemize}
\begin{table}[h!]
	\centering
	\begin{tabular}{ll}
		\hline
		Windows & Linux, OSX\\
		\hline 
		Win32 Disk Imager & dd\\
		Advanced IP scanner & nmap\\
		Putty & ssh\\
		WinSCP & sftp\\
		\hline
	\end{tabular}
	\caption{RPi installation and setup utilities}
	\label{tab:RPi_utilities} 
\end{table}
See Table \ref{tab:RPi_utilities} for a list of the equivalent software for OSX and Linux.
\subsubsection{Write image to SD card}
Since the .iso file is raw, it needs to be written to the SD card in way that makes it bootable. Win32 Disk Imager does this.
\begin{figure}[htb!]
	\centering \includegraphics[scale=0.45]{fig/Rpi_DiskImager} \caption{Disk Imager}
	\label{fig: Disk Imager} 
\end{figure}
Run the program as administrator. Select the correct image file and device, as in Figure \ref{fig: Disk Imager}. Make sure that you have selected the correct drive before you push \texttt{Write}. Once the write is complete, insert the SD card in the RPi and boot.
\subsubsection{Terminal access}\label{subsec: Terminal access}
RPi can be accessed through the network, i.e. without having to directly connect a monitor and keyboard.
\begin{figure}[htb!]
	\centering \includegraphics[scale=0.45]{fig/advancedIPscanner} \caption{Advanced IP Scanner}
	\label{fig: Advanced IP Scanner} 
\end{figure} 
At first boot, the RPi by default waits to be assigned an IP address by DHCP. If this address is not known, scan the network with Advanced IP Scanner. It is advicible to sort the results by manufacturer since it is fixed (\textit{Raspberry Pi Foundation}). The name is typically \emph{raspberrypi}. See Figure \ref{fig: Advanced IP Scanner}.
\begin{figure}[htb!]
	\centering \includegraphics[scale=0.45]{fig/Rpi_remote_access1} \caption{Putty settings}
	\label{fig: Putty settings} 
\end{figure}
Once the IP is known, it is specified in the Putty settings, as in Figure \ref{fig: Putty settings}, and a connection can be opened. 
\begin{figure}[htb!]
	\centering \includegraphics[scale=0.45]{fig/Rpi_remote_access2} \caption{SSH connection}
	\label{fig: SSH connection} 
\end{figure}
The default login is \texttt{pi}, and the default password \texttt{raspberry}. Figure \ref{fig: SSH connection} shows the terminal output on first login.
\subsubsection{Finalize configuration}
\begin{figure}[htb!]
	\centering 
	\includegraphics[scale=0.45]{fig/Rpi_finalize_install} 
	\caption{RPi configuration tool}
	\label{fig: RPi configuration tool} 
\end{figure}
Enter the\begin{verbatim}sudo raspi-config\end{verbatim}command to start the RPi Software Configuration Tool, as in Figure \ref{fig: RPi configuration tool}. Use the menu to apply the following
\begin{enumerate}
	\item Update configuration tool: 8 Advanced Options \textgreater  A9 Update
	\item Change password: 2 Change User Password
	\item Expand filesystem: 1 Expand Filesystem \textgreater  Finish
\end{enumerate}
Exit the configuration tool and select \texttt{Yes} for reboot. Reconnect through Putty.
Finally, update the repository package lists and upgrade all packages currently installed on the RPi:\begin{verbatim}sudo apt-get update 
sudo apt-get upgrade -y\end{verbatim}This process took approximately 10 minutes on a 90 Mbps internet connection. 
\subsubsection{Transfer files to RPi from computer}\label{subsec: Transfer files to RPi}
WinSCP can be used to transfer files to the RPi. This is useful for instance when transferring code, or when the RPi is not directly connected to the internet.
\subsubsection{Set fixed IP address}
When the RPi is connected directly to the cRIO or computer, a fixed IP is necessary since there is no DHCP server in that network. During most of this setup, however, it is preferable to keep the default DHCP assigned IP setting. To set a fixed IP
\begin{enumerate}
	\item Open the network interface configuration information file for editing\begin{verbatim}sudo nano /etc/network/interfaces\end{verbatim}
	\item Alter the eth0 settings from \texttt{dhcp} to \texttt{static} and
	add address and netmask as\begin{verbatim}auto eth0
	iface eth0 inet static
	address 192.168.1.22
	netmask 255.255.255.0\end{verbatim}
	\item Save the changes by the key combination \texttt{Ctrl+X}.
\end{enumerate}
The new IP is applied on the next reboot.

\subsection{Sixaxis installation and configuration}
This section describes how to install and configure the Sixaxis gamepad for Bluetooth connection to the RPi, and how to add a server for sending joystick signals to the cRIO.
\subsubsection{Download and install bluetooth support}
BlueZ is the official Linux Bluetooth stack. It provides support for core Bluetooth layers and protocols. To download and install, type
\begin{verbatim}sudo apt-get install bluez-utils bluez-compat bluez-hcidump
libusb-dev libbluetooth-dev joystick checkinstall -y\end{verbatim}
The process takes a few minutes.
\begin{figure}[htb!]
	\centering \includegraphics[scale=0.45]{fig/rpi_hciconfig} \caption{Bluetooth configuration tool}
	\label{fig: RPi hciconfig} 
\end{figure}
To confirm the installation, use the \texttt{hciconfig} command to print name and basic information about Bluetooth devices installed in the system. The output should include \texttt{UP RUNNING PSCAN}, as in Figure \ref{fig: RPi hciconfig}. If instead it says \texttt{DOWN}, some error har occured. Most experienced errors were due to typos.
\subsubsection{Bluetooth pairing}\label{par: Bluetooth-pairing}
Sixaxis does not support the standard Bluetooth paring prcedure, instead, pairing is done over USB. The \texttt{sixpair} command-line utility\footnote{by Pabr Technologies, www.pabr.org} searches USB buses for Sixaxis devices and tells them to connect to a new Bluetooth master.

Download and compile the program by the following commands:
\begin{verbatim}wget http://www.pabr.org/sixlinux/sixpair.c
gcc -o sixpair sixpair.c -lusb\end{verbatim}
Connect the Sixaxis by USB before running the paring utility
\begin{verbatim}sudo ./sixpair\end{verbatim}
The output should be similar to
\begin{verbatim}Current Bluetooth master: 00:02:72:BF:BC:8F
Setting master bd_addr to: 00:02:72:BF:BC:8F\end{verbatim}
The addresses at the end of each line will only be the same if you have already paired the Sixaxis with the Bluetooth dongle. First time they will be different. The Sixaxis USB cable may now be disconnected.

\subsubsection{Joystick manager system service}
\texttt{QtSixA}\footnote{the Sixaxis Joystick Manager by falkTX, qtsixa.sourceforge.net} reads the Sixaxis signals and makes them available to other programs. This program needs to run automatically whenever the RPi is booted.

To download the program, type
\begin{verbatim}wget http://sourceforge.net/projects/qtsixa/files/QtSixA%201.5.1/QtSixA-1.5.1-src.tar.gz\end{verbatim}
To install, type

\begin{verbatim}tar xfvz QtSixA-1.5.1-src.tar.gz
cd QtSixA-1.5.1/sixad
make
sudo mkdir -p /var/lib/sixad/profiles
sudo checkinstall -y\end{verbatim}
Update the system service list with sixad driver and reboot
\begin{verbatim}sudo update-rc.d sixad defaults
sudo reboot\end{verbatim}
To test the program, turn on the Sixaxis (round PS button in the middle) and start the test program
\begin{verbatim}sudo jstest /dev/input/js0\end{verbatim}
The terminal should now fill up with numbers that change as you move the analogue sticks and press the buttons on the Sixaxis. Exit the program by the key combination \texttt{Ctrl+C}.

\subsubsection{Joystick signal server}
A server must run to make joystick signals available over the RPi ethernet port. This should also start whenever the RPi is booted.

Transfer the source file \texttt{jscont.c} to the RPi (see Section
\ref{subsec: Transfer files to RPi}), then compile:

\begin{verbatim}g++ -o jscont jscont.c\end{verbatim}

To verify that the program runs correctly, turn off (hold PS3 button
for about 10 seconds) the previously paired Sixaxis and start the
program

\begin{verbatim}./jscont\end{verbatim}

\begin{figure}[h!]
	\centering \includegraphics[scale=0.45]{fig/RPi_jscont} \caption{Joystick signal server test}
	\label{fig: RPi joystick server} 
\end{figure}

The program should then wait until you turn on the Sixaxis before
giving output simular to Figure \ref{fig: RPi joystick server}. To
exit the server use the key combination \texttt{Ctrl+C}.

Next, disable login at start-up in the bootup service description
\texttt{inittab}:
\begin{enumerate}
	\item Open the file for editing \begin{verbatim}sudo nano /etc/inittab\end{verbatim}
	\item Change the line that reads \begin{verbatim}1:2345:respawn:/sbin/getty --noclear 38400 tty1\end{verbatim}
	by adding \texttt{-{}-autologin pi} to get \begin{verbatim}1:2345:respawn:/sbin/getty --autologin pi --noclear 38400 tty1\end{verbatim}
	\textbf{Warning:} Typos here may result consequences hard to correct.
	\item Save and exit the changes by the key combination \texttt{Ctrl+X}.
\end{enumerate}
Finally, add \texttt{jscont} to the login execution file:
\begin{enumerate}
	\item Open the file for editing \begin{verbatim}sudo nano /home/pi/.bashrc\end{verbatim}
	\item At the very end of the file, add \begin{verbatim}sudo ./jscont\end{verbatim}
	\item Save the changes by the key combination \texttt{Ctrl+X}.
\end{enumerate}
RPi should now be sending joystick signals at start-up.
%%=========================================
\bibliographystyle{elsarticle-harv}
\addcontentsline{toc}{chapter}{\bibname}
\bibliography{refs}  
%%=========================================

\end{document}
